%==============================================================================
% Paper G -- skeleton (bullets, to be polished into prose)
% Target venue: systems / methodology (e.g., a benchmarking or virtualization
%   workshop). This paper documents the SYNTHETIC GUEST WORKLOADS used by the
%   rest of the portfolio: their design rationale, the safety argument, the
%   tunable knobs, and how the two families map onto the downstream signal.
%   It deliberately does NOT re-state Paper C's classification contribution --
%   it only characterizes the workloads that Paper C consumes.
%
% Build: pdflatex skeleton.tex   (needs IEEEtran; if unavailable, change the
%   \documentclass line below to \documentclass[11pt]{article} and delete the
%   \IEEEauthorblock* macros.)
%==============================================================================
\documentclass[conference]{IEEEtran}
\usepackage{amsmath,booktabs,siunitx,url}

\title{Reversible-XOR Memory Workloads: A Synthetic Sandbox\\
Taxonomy for Introspection Benchmarks}

\author{\IEEEauthorblockN{TODO Author list}
\IEEEauthorblockA{TODO Affiliation\\
\texttt{TODO@email}}}

\begin{document}
\maketitle

%------------------------------------------------------------------
\begin{abstract}
% TODO: polish into 150-200 words of prose. Bullet stub:
\begin{itemize}
  \item Problem: exercising whole-VM memory-change capture (Paper A's Active
        Page Fraction, APF) needs controllable guest workloads that drive
        memory at known intensities and rhythms, yet must be SAFE to run
        repeatedly in a research sandbox -- no real malware, no network, no
        persistence.
  \item Idea: an 11-workload taxonomy split into two families -- \emph{phasic}
        (rhythmic active/idle) and \emph{steady} (uniform) -- all built on
        \emph{reversible XOR} memory operations over a bounded working-set
        buffer, so every run is fully reversible and self-contained.
  \item Knobs: each workload exposes a small set of tunable parameters --
        working-set size (\texttt{--working-set-mb}), intensity, and (for
        phasic) phase period/cadence -- documented as a workload~$\times$~knob
        table.
  \item Signal coupling: phasic workloads produce a rhythmic APF trajectory
        (cepstral-SNR is the defining metric; phasic SNR floor \num{4.78}),
        steady workloads produce a near-constant APF (low coefficient of
        variation). Both families feed the downstream tuning (Paper B),
        classification (Paper C), and segmentation (Paper D).
  \item Honesty: one workload labeled \emph{steady}
        (\texttt{mem\_pagefault\_density\_v2}) does \emph{not} behave as steady
        under segmentation analysis -- a documented taxonomy caveat, not a
        result we hide.
\end{itemize}
\end{abstract}

%------------------------------------------------------------------
\section{Introduction}
% Proposed introduction (bullets):
\begin{itemize}
  \item Hypervisor-external memory introspection (Paper A) observes a guest's
        memory churn without an in-guest agent. To study and tune such a
        capture pipeline, one needs guest workloads whose memory behavior is
        \emph{known}, \emph{controllable}, and \emph{repeatable}.
  \item Real malware (the motivating threat model:
        ransomware-like encryption bursts) is unsafe to run repeatedly on a
        shared research host and is not reproducible. We therefore build
        \emph{synthetic} guest workloads that imitate the relevant
        \emph{memory-access shapes} without any harmful effect.
        % HONESTY: these are imitations of memory-access SHAPES, not validated
        % behavioral replicas of real ransomware. The "ransom" naming is by
        % analogy (discover -> read -> XOR -> write -> rename), not a claim of
        % fidelity to any specific malware family. \cite{TODO_ransomware_survey}
  \item \textbf{Design constraints} (the safety argument, Sec.~\ref{sec:safety}):
        (a) memory mutation uses \emph{reversible XOR} over a working-set
        buffer; (b) no network; (c) no persistence beyond a validated sandbox
        directory; (d) bounded working-set (the binding knob is
        \texttt{--working-set-mb}, e.g.\ the sweep goes up to \SI{1024}{\mebi\byte}
        within a \SI{1}{\gibi\byte} guest).
  \item \textbf{Two families.} We organize the catalog into \emph{phasic}
        (rhythmic active/idle, e.g.\ bursty or paced file processing) and
        \emph{steady} (uniform, continuous memory churn). The split is not
        cosmetic: it predicts the \emph{shape} of the downstream APF signal
        (Sec.~\ref{sec:coupling}) and is the axis the classifier recovers
        (Paper C).
  \item \textbf{This paper} documents the workload-design rationale, the safety
        model, the tunable knobs, and the family$\to$signal mapping. It does
        \emph{not} re-derive the classification result (Paper C) or the
        segmentation method (Paper D); those are the \emph{downstream uses} of
        these workloads and are forward-referenced.
  \item \textbf{Contributions} (expanded in Sec.~\ref{sec:contrib}):
    \begin{itemize}
      \item C1: an 11-workload, two-family synthetic taxonomy for memory-change
            introspection benchmarks;
      \item C2: a safety argument (reversible XOR, no network, no persistence,
            bounded working-set) making the workloads research-sandbox-safe;
      \item C3: a tunable-knob specification and the family$\to$signal mapping,
            including an honest caveat where a designed-steady workload
            misbehaves (Sec.~\ref{sec:honesty}).
    \end{itemize}
\end{itemize}

%------------------------------------------------------------------
\section{Background and Related Work}
\begin{itemize}
  \item Synthetic benchmarks / microbenchmarks for memory systems: the tradition
        of driving a memory subsystem with parameterized synthetic loads rather
        than opaque applications. TODO: cite a representative memory
        microbenchmark suite. \cite{TODO_working_set_estimation}
  \item Ransomware behavior and detection surveys: motivate the phasic
        (encryption-burst / file-walk) shapes without running real malware.
        \cite{TODO_ransomware_survey}
        % HONESTY: we cite the threat model for MOTIVATION only; we make no
        % detection claim here (that is Paper C). No real malware is executed.
  \item Periodicity / cepstral analysis: support the cepstral-SNR metric used to
        characterize phasic rhythm. \cite{TODO_cepstrum}
  \item Working-set estimation / dirty-page tracking: context for what "memory
        activity" means at the hypervisor boundary; APF (Paper A) is the
        content-derived signal these workloads are designed to move.
        \cite{TODO_working_set_estimation}
  \item TODO: position against existing safe-execution / sandbox harnesses for
        malware analysis (we differ by being SYNTHETIC and reversible rather
        than containing/observing real samples).
\end{itemize}

%------------------------------------------------------------------
\section{Workload Taxonomy}
\label{sec:taxonomy}
\begin{itemize}
  \item \textbf{Source of truth.} The catalog and per-workload one-line
        behaviors are documented in \texttt{workload\_scope\_rationale.md}; the
        v3 campaign that instantiated them is \texttt{PLAN02\_V3\_RUNBOOK.md}.
        All binaries live under \texttt{VM\_executables\_phase2/} and are
        reversible XOR synthetic probes (no real malware).
  \item \textbf{Two families.} \emph{Phasic} = rhythmic active/idle;
        \emph{steady} = uniform/continuous. The runbook also groups the same
        binaries into design \emph{groups} (security-like, memory,
        app-realistic, methodology); the family axis (phasic/steady) is the
        \emph{signal-shape} axis used downstream and is what we report here.
  \item \textbf{The 11 workloads} (confirmed against the sources; family =
        phasic/steady):
    \begin{itemize}
      \item PHASIC -- \texttt{sandbox\_ransom\_batched}: all files in one batch
            through 5 phases, then exits (one short burst). % bursty probe
      \item PHASIC -- \texttt{sandbox\_ransom\_seq}: per-file loop
            (stat $\to$ read $\to$ XOR $\to$ write $\to$ rename); sustained,
            many phase boundaries.
      \item PHASIC -- \texttt{sandbox\_ransom\_slowburn}: paced, one file every
            \texttt{--interval-s} seconds (sustained + periodic cadence).
      \item PHASIC -- \texttt{sandbox\_ransom\_selective}: extension-filtered
            discovery then process (burst + discovery pre-phase).
      \item PHASIC -- \texttt{sandbox\_scanner\_metadata}: metadata-only scan
            (low, flat -- near-idle control; multi-pass produces a distinctive
            cepstral peak per Paper C).
      \item STEADY -- \texttt{mem\_workingset\_sweep\_v2}: cyclic working-set
            churn (steady, continuous activity). The canonical steady probe.
      \item STEADY -- \texttt{mem\_mmap\_traversal\_v2}: mmap region traversal.
      \item STEADY -- \texttt{mem\_pagefault\_density\_v2}: page-fault-density
            driver. % HONESTY: labeled steady but misbehaves -- see Sec. honesty.
      \item STEADY -- \texttt{mem\_rmw\_intensity\_v2}: read-modify-write
            intensity sweep.
      \item STEADY -- \texttt{mem\_writemag\_sweep\_v2}: write-magnitude sweep.
      \item STEADY -- \texttt{app\_hashtable\_intensive\_v2}: hashtable
            build $\to$ probe (two-phase / bimodal trajectory). Labeled steady,
            scored with CV primary plus incidental F1; a dedicated bimodal
            metric is deferred.
            % HONESTY: app_hashtable is BIMODAL (build->probe), not strictly
            % uniform; it is filed under "steady" for metric dispatch, with a
            % bimodal metric deferred (per PLAN02_V3_RUNBOOK notes).
    \end{itemize}
  % All 11 confirmed: 5 phasic + 6 steady (cross-checked against
  % workload_scope_rationale.md S.2 and workload_classification_results.html S.3).
  \item TODO: figure -- one example APF trajectory per family (a phasic
        burst train vs.\ a steady plateau).
  \item TODO: table -- the 11 workloads $\times$ knobs (working-set-mb,
        intensity/mode, phase period/cadence, family). Skeleton below
        (Table~\ref{tab:knobs}); fill remaining cells from the binaries'
        \texttt{--help} -- mark any unconfirmed knob as TODO rather than
        inventing a value.
\end{itemize}

% TODO: table -- 11 workloads x knobs. Values shown are the v3-campaign
% instantiations from PLAN02_V3_RUNBOOK Step 3 (NOT the full knob ranges).
% Empty/"--" cells need confirmation from each binary's --help.
\begin{table}[t]
  \centering
  \caption{Workloads $\times$ tunable knobs (v3 instantiation; see
           \texttt{PLAN02\_V3\_RUNBOOK.md} Step~3). ws-mb $=$
           \texttt{--working-set-mb}. Blank $=$ % TODO: confirm from --help.}
  \label{tab:knobs}
  \begin{tabular}{@{}llll@{}}
    \toprule
    Workload & Family & Key knob(s), v3 value & Notes \\
    \midrule
    ransom\_batched   & phasic & \texttt{--files 12000} & 1\,MiB files; one burst \\
    ransom\_seq       & phasic & \texttt{--files 4000}  & 1\,MiB files; sustained \\
    ransom\_slowburn  & phasic & \texttt{--files 200 --interval-s 3} & paced \\
    ransom\_selective & phasic & \texttt{--files 1500}  & 1\,MiB files; filtered \\
    scanner\_metadata & phasic & \texttt{--files 5000 --passes 40} & metadata only \\
    workingset\_sweep & steady & \texttt{--working-set-mb 256} & \texttt{--stride 4096} \\
    mmap\_traversal   & steady & \texttt{--variant rmw --file-size-mb 256} & \\
    pagefault\_density& steady & \texttt{--variant mixed --working-set-mb 256} & see honesty \\
    rmw\_intensity    & steady & \texttt{--mode rmw --working-set-mb 256} & \texttt{--stride 4096} \\
    writemag\_sweep   & steady & \texttt{--working-set-mb 256 --bytes-per-page 64} & \\
    hashtable\_intensive & steady & \texttt{--capacity-pow2 24} & bimodal; CV primary \\
    \bottomrule
  \end{tabular}
  % HONESTY: the rationale doc states the working-set knob ranges UP TO
  % --working-set-mb 1024 (full 1 GiB guest); the v3 campaign instantiated the
  % steady probes at 256. The table reports the v3 value, not the max range.
\end{table}

%------------------------------------------------------------------
\section{Safety Model}
\label{sec:safety}
\begin{itemize}
  \item \textbf{Reversibility (XOR).} Memory mutation is XOR over a working-set
        buffer; applying the same key twice restores the original bytes, so a
        run leaves no irreversible state change in the touched buffer. This is
        the core safety property: the workloads \emph{churn} memory without
        \emph{destroying} data.
        % TODO: confirm whether the file-processing ("ransom") probes XOR the
        % sandbox files in place and rename, and whether that is reversed at
        % exit, vs. operating on a throwaway working-set copy. State precisely.
  \item \textbf{No network.} No workload opens a network connection; all
        activity is local memory/file churn. (Asserted by
        \texttt{PLAN02\_V3\_RUNBOOK.md}.)
  \item \textbf{No persistence / validated sandbox.} File-touching probes
        operate only inside a validated sandbox directory; nothing escapes it.
        (Asserted by the runbook: "validated sandbox dir.")
  \item \textbf{Bounded working-set.} The memory footprint is capped by
        \texttt{--working-set-mb} (v3: \SI{256}{\mebi\byte}; rationale documents
        the knob up to \SI{1024}{\mebi\byte} within the \SI{1}{\gibi\byte}
        guest), so a workload cannot exhaust host or guest memory.
  \item \textbf{Not real malware.} Every binary is a synthetic probe; the
        "ransom" naming denotes an imitated \emph{access shape}
        (discover/read/XOR/write/rename), not a malware sample.
        % HONESTY: safety here is by CONSTRUCTION + assertion in the runbook,
        % not an independent audit. Before publication: state how reversibility
        % and sandbox-confinement were VERIFIED (test target? checksum
        % round-trip?), not just intended. The runbook references `make smoke`.
  \item TODO: figure/box -- the safety checklist (reversible / no-net /
        no-persist / bounded) as a one-glance contract.
\end{itemize}

%------------------------------------------------------------------
\section{Contributions in Detail}
\label{sec:contrib}
\begin{itemize}
  \item \textbf{C1 -- taxonomy}: 11 synthetic workloads in two families
        (5 phasic, 6 steady) spanning bursty, paced, sustained, and uniform
        memory-access shapes; the family axis is the signal-shape axis the rest
        of the portfolio uses.
  \item \textbf{C2 -- safety}: a reversible-XOR, no-network, no-persistence,
        bounded-working-set construction that makes repeated whole-VM capture
        safe on a shared research host.
  \item \textbf{C3 -- knobs + coupling}: a workload$\times$knob specification
        and the empirical family$\to$signal mapping (phasic $\Rightarrow$
        rhythmic APF, cepstral SNR $\ge \num{4.78}$; steady $\Rightarrow$
        near-constant APF, low CV), with one honest exception
        (Sec.~\ref{sec:honesty}).
\end{itemize}

%------------------------------------------------------------------
\section{Family-to-Signal Coupling}
\label{sec:coupling}
\begin{itemize}
  \item \textbf{The signal.} APF (Paper A) is the fraction of \SI{4}{\kibi\byte}
        guest pages that changed between consecutive \SI{1}{\gibi\byte}
        snapshots ($N = 262{,}144$ pages). A workload is "visible" through how
        it moves APF over time.
  \item \textbf{Phasic $\Rightarrow$ rhythmic APF.} Active/idle workloads
        produce a periodic APF trajectory; the defining metric is the cepstral
        signal-to-noise ratio (cepstral-SNR). Phasic SNR floor observed:
        \num{4.78} (per-cell phasic SNR in the range $\approx$\num{4.7}--\num{6.4}
        \si{\decibel} per Paper C). \cite{TODO_cepstrum}
  \item \textbf{Steady $\Rightarrow$ near-constant APF.} Uniform workloads
        produce a flat APF; the defining metric is the coefficient of variation
        (CV), low for steady cells ($\approx$\num{0.06}--\num{0.33} when
        populated, per Paper C).
  \item \textbf{Downstream use (forward references).} These workloads are the
        common substrate for: Paper B (single-interval gated tuning -- the
        cadence-floor argument), Paper C (phasic-vs-steady classification, the
        DOWNSTREAM use of this taxonomy; we do NOT duplicate its 100\%-binary
        result here), and Paper D (CUSUM segmentation, where the steady gate is
        applied -- and where one workload fails, Sec.~\ref{sec:honesty}).
  \item TODO: table/figure -- per-family summary stat (phasic: cepstral-SNR;
        steady: CV) aggregated over the v3 cells, with the pagefault row
        flagged. Pull from the v3 metrics; do not duplicate Paper C's confusion
        matrix.
\end{itemize}

%------------------------------------------------------------------
\section{The v3 Campaign (Instantiation)}
\label{sec:campaign}
\begin{itemize}
  \item \textbf{Design.} Campaign v3 instantiated all 11 workloads at a single
        sampling interval (\SI{500}{\milli\second}; the snapshot cadence floor
        makes finer intervals moot, per Paper A/B) across 3 durations
        $\{120, 300, 600\}\,\si{\second}$.
  \item \textbf{Cells.} \textbf{132 cells} $= 11$ workloads $\times 3$ durations
        $\times 4$ replicates; \textbf{128 ok}.
        % HONESTY: PLAN02_V3_RUNBOOK Step 3 runs --replicates 2 -> 66 production
        % cells per bundle. The 4 replicates / 132 cells come from COMBINING two
        % bundles (v3_full + v3_ext, 2 replicates each), as
        % workload_classification_results.html S.02-03 documents. Report 132/4
        % per the verified ground-truth; note the two-bundle merge.
  \item \textbf{Per-workload yield (usable at the analyzer setting W=8,H=4).}
        Most workloads yield 9--12 usable cells; 14 of 132 are dropped for short
        trajectories or numerically unstable cepstra (per Paper C's dataset
        table). This paper uses those yields only to characterize the workloads;
        the classifier itself is Paper C.
  \item TODO: confirm the 600\,s ceiling rationale -- the binaries reject
        \texttt{--duration > 600} (runbook "Notes + caveats"); a longer sweep
        must clamp in \texttt{\_augment\_workload\_command}.
\end{itemize}

%------------------------------------------------------------------
\section{Honesty: A Designed-Steady Workload That Misbehaves}
\label{sec:honesty}
\begin{itemize}
  \item \texttt{mem\_pagefault\_density\_v2} is \emph{labeled} a steady-family
        workload and \emph{passes} the Plan~03 stationarity/coverage gates, but
        it \emph{fails} the Plan~04 steady segmentation gate (Paper D): its APF
        is not truly stationary.
  \item \textbf{Evidence (by duration).} Mean APF swings with capture duration:
        $\approx$\num{0.016} at \SI{120}{\second}, $\approx$\num{0.497} at
        \SI{300}{\second}, $\approx$\num{0.257} at \SI{600}{\second}. The
        steady-metric CV is unstable too: \texttt{cv\_workingset} $\approx$
        \num{0.88} at \SI{120}{\second} vs.\ $\approx$\num{0.10} at
        \SI{300}{\second}.
        % HONESTY (core caveat of this paper): a workload we DESIGNED to be
        % steady does not behave as steady under segmentation. We surface this
        % rather than relabel it. Likely cause to investigate: fault-density
        % driver's footprint interacts with duration / first-touch population
        % of the working set. Do NOT claim a mechanism we have not verified.
  \item \textbf{Taxonomy implication.} The phasic/steady label is a
        \emph{design intent}, validated downstream -- and it is not infallible.
        We report one labeled-steady workload that the steady gate rejects, as
        a caveat on the taxonomy's completeness, not as a hidden failure.
  \item Note the related, separately-flagged case:
        \texttt{app\_hashtable\_intensive\_v2} is bimodal (build$\to$probe) yet
        filed as steady for metric dispatch; a dedicated bimodal metric is
        deferred. (Distinct from the pagefault non-stationarity above.)
  \item TODO: table -- pagefault APF/CV per duration (120/300/600), side by side
        with a well-behaved steady probe (e.g.\ workingset\_sweep), to make the
        misbehavior legible.
\end{itemize}

%------------------------------------------------------------------
\section{Limitations}
\begin{itemize}
  \item \textbf{Imitation, not fidelity.} The "ransom" probes imitate
        memory-access \emph{shapes}; they are not validated replicas of any real
        ransomware family, and no detection claim is made here (that is
        Paper C). \cite{TODO_ransomware_survey}
  \item \textbf{Clean exemplars, not a population.} The 11 workloads are
        deliberately clean single-mode (or bimodal) exemplars. Mixed
        phasic+steady behavior (e.g.\ a database with periodic checkpoints) is
        not represented (echoing Paper C's "not a population sample" caveat).
  \item \textbf{Label is intent, not guarantee.} As Sec.~\ref{sec:honesty}
        shows, a designed-steady workload can fail a steady gate; the taxonomy
        is a useful organizing axis, not a proof of behavior.
  \item \textbf{Safety by construction, not audit.} Reversibility and
        sandbox-confinement are asserted/constructed and exercised by a smoke
        target, not independently audited; precise verification evidence is
        TODO before publication.
  \item \textbf{Single host/guest, single day.} All replicates share one host,
        one guest, one capture day (per Paper C's generalization caveat); the
        knob$\to$signal mapping is not yet shown to be host-invariant.
\end{itemize}

%------------------------------------------------------------------
\section{Conclusion}
\begin{itemize}
  \item A small, SAFE, reversible-XOR synthetic taxonomy -- 11 workloads in two
        families (phasic, steady) -- is enough to exercise whole-VM
        memory-change capture across bursty, paced, sustained, and uniform
        access shapes, with bounded working-sets and no network or persistence.
  \item The family axis predicts the APF signal shape (rhythmic vs.\ constant)
        and is the substrate the rest of the portfolio tunes (Paper B),
        classifies (Paper C), and segments (Paper D) on.
  \item We keep the taxonomy honest: one labeled-steady workload
        (\texttt{mem\_pagefault\_density\_v2}) is not stationary under
        segmentation, a caveat we document rather than conceal.
\end{itemize}

\bibliographystyle{IEEEtran}
\bibliography{../refs}
\end{document}
